\section{自动微分：数值与梯度的分别实现}
\label{sec:kl-autodiff}
前面的统一框架估计已经指定的积分。本节检查标量程序的反传究竟对应哪个积分。若一个程序同时用于 KL 日志与训练，审计对象应记录为 forward 数值和 backward 向量组成的有序对。
本节涉及随机梯度的期望与偏差时默认 $\rho=p$；固定索引的逐样本求导规则本身不依赖抽样分布。

\subsection{一个函数的两种求导路径}

\begin{theorem}[固定样本反传与采样测度补偿]
\label{thm:kl-autodiff-transform}
令 $\phi$ 是不另依赖参数 $\theta$ 的可微实函数，$A\sim p$，反传时固定样本索引。定义
\[
 F_A=\phi(r_A),\qquad
 \kappa_A=\frac{p(A)}{\sg(p(A))},\qquad
 L_A=\kappa_A\phi(r_A).
\]
两者的 forward 值逐样本相同，固定样本导数分别为
\[
 \nabla F_A=\phi'(r_A)s_A,\qquad
 \nabla L_A=\{\phi(r_A)+\phi'(r_A)\}s_A.
\]
后者的期望等于 $\nabla I_p(\phi(r))$。若仅将系数停止梯度，$\kappa_A\sg(\phi(r_A))$ 的导数则为 $\phi(r_A)s_A$。
\end{theorem}

\begin{proof}
固定索引下 $\nabla r_A=s_A$，链式法则给出第一式。forward 时 $\kappa_A=1$；分母停止梯度，因而 $\nabla\kappa_A=s_A$。乘积法则给出第二式。再由式~\eqref{eq:kl-measure-derivative}，
\[
 \nabla I_p(\phi(r))
 =I_p\bigl(\phi'(r)s+\phi(r)s\bigr),
\]
等于第二种导数的期望。停止整个 $\phi$ 时只保留 $\nabla\kappa_A$ 一项，得到最后的规则。
\end{proof}

这个变换把 K1--K4 的差异归到同一个问题：给定 forward 函数后，程序保留了哪些乘积法则项。数值无偏是 $I_p(\phi(r))=K$，裸反传无偏则是 $I_p(\phi'(r)s)=g$，两者是不同等式。

\begin{center}
\begin{tabular}{lll}
\toprule
程序 & forward 值 & 固定样本 backward\\
\midrule
K1 & $r_A$ & $s_A$\\
K2 & $r_A^2/2$ & $r_As_A$\\
K3 & $e^{-r_A}+r_A-1$ & $(1-e^{-r_A})s_A$\\
K4 & $r_A$ & $r_As_A$\\
K4$^{\mathrm{live}}$ & $r_A$ & $(r_A+1)s_A$\\
K3$^{++}$ & $e^{-r_A}+r_A-1$ & $r_As_A$\\
\bottomrule
\end{tabular}
\end{center}
数值方差按第二列计算，梯度方差按第三列计算。K5 的目标是正向 KL，随后单独代入同一求导规则。

\subsection{K4 与固定样本的梯度替换}
K1 与 K3 的数值都无偏；K2 是以二阶矩换取局部稳定性的有偏 proxy。若需要同时让数值和自动微分梯度对应同一个 $K$，使用停止 log-ratio 系数的 K4：
\[
 \kappa_A=\frac{p(A)}{\operatorname{sg}(p(A))},\qquad
 \widehat K_4=\kappa_A\sg(r_A).
\]
数值上 $\kappa_A=1$，所以 $\E[\widehat K_4]=K$；反传时
\[
 \nabla_\theta\widehat K_4=r_As_A,
\]
因此
\[
 \E_p[\nabla_\theta\widehat K_4]=\nabla_\theta K=g.
\]
K4 与 K1 的 forward 值逐样本相同，故 K4 的数值偏差为零，数值方差也为 $\Var_p(r)$。stop-gradient 只改变求导规则。另定义本项目单 token 测试采用的 live-log-ratio 变体
\[
 \widehat K_4^{\mathrm{live}}=\kappa_A r_A,\qquad
 \nabla\widehat K_4^{\mathrm{live}}=(r_A+1)s_A.
\]
它的数值仍与 K4 相同，梯度期望也相同，但梯度方差不同。后文始终保留上标 $\mathrm{live}$，不将两种反传路径共用一个名称。

\begin{theorem}[K4 的数值与梯度无偏性]
\label{thm:k4}
在本章正支持与固定参考假设下，$\widehat K_4$ 和 $\widehat K_4^{\mathrm{live}}$ 的数值均对 $K$ 无偏，自动微分梯度均对 $g$ 无偏。
\end{theorem}
\begin{proof}
两者的数值为 $r_A$，期望是 $\sum_vp(v)r_v=K$。两者的梯度分别为 $r_As_A$ 与 $(r_A+1)s_A$，期望之差为 $\sum_vp(v)s_v=0$，因此都等于 $g$。
\end{proof}

\subsection{K3$^{++}$ 与正向 KL 的 K5}

EMA-PG 的公开实现使用上述停止系数 K4，并定义 K3$^{++}$ 与 K5；以下公式先去掉 clipping，单独分析数学对象。\footnote{代码版本固定为 EMA-PG \texttt{f9b9776}，定义见该版本的 \href{https://github.com/LunjunZhang/ema-pg/blob/f9b977632ad948ce935432a517c22291bfaba562/search/verl/trainer/ppo/core_algos.py}{\texttt{core\_algos.py}}。}
\[
 \widehat K_{3^{++}}=\kappa_A(e^{-r_A}+r_A-1).
\]
它的数值与 K3 相同，因此对反向 $K$ 无偏，数值方差也与 K3 相同。其反传为
\[
 \nabla\widehat K_{3^{++}}
 =\{e^{-r_A}+r_A-1+1-e^{-r_A}\}s_A=r_As_A.
\]
因此它与 K4 的 backward 逐样本相同，梯度偏差为零，梯度方差也完全相同；数值方差的差异不能解释成梯度方差的差异。

K5 对应正向目标 $K^{\mathrm F}=\KL(q\Vert p)$，定义为
\[
 \widehat K_5=-\sg(e^{-r_A})r_A+\log p(A)-\sg(\log p(A)).
\]
其 forward 值为 $-e^{-r_A}r_A$。在 $A\sim p$ 下，
\[
 \E_p[\widehat K_5]=-\sum_vq(v)r_v=K^{\mathrm F},\qquad
 \Var(\widehat K_5)=\E_p[e^{-2r_A}r_A^2]-(K^{\mathrm F})^2.
\]
它相对本章默认反向目标 $K$ 的数值偏差是 $K^{\mathrm F}-K$，不能与 K1--K4 按同一个无偏目标直接并列。反传为
\[
 \nabla\widehat K_5=(1-e^{-r_A})s_A,\qquad
 \E_p[\nabla\widehat K_5]=-\sum_vq(v)s_v=\nabla K^{\mathrm F}.
\]
其梯度方差为
\[
 \E_p[(1-e^{-r_A})^2\lVert s_A\rVert^2]
 -\left\lVert\sum_vq(v)s_v\right\rVert^2.
\]
这里正向目标只改变 KL 的积分方向，student 与 score 的定义保持不变。

若数值稳定化把 K3 的函数 $\phi(r)=e^{-r}+r-1$ 改成分段可导的 $\widetilde\phi(r)$，K3$^{++}$ 的数值偏差和梯度偏差分别变成
\[
 \E_p[\widetilde\phi(r_A)]-K,\qquad
 \E_p[(\widetilde\phi(r_A)+\widetilde\phi'(r_A)-r_A)s_A].
\]
公开代码对 log-ratio 与输出均有 clamp，因此无 clipping 恒等式的适用条件是相关 clamp 未激活，或相关偏差项经额外核验为零。

\subsection{数值 estimator 的自动微分梯度}

把已经抽到的动作视为固定索引，对 K1--K4 直接反传，得到
\[
 G_1=s_A,\quad G_2=r_As_A,\quad
 G_3=(1-e^{-r_A})s_A,\quad G_4=r_As_A.
\]
令 $u_q=\E_q[s_A]$，则这些随机梯度相对于 $g$ 的偏差为
\[
 \operatorname{Bias}(G_1)=-g,\quad
 \operatorname{Bias}(G_2)=0,\quad
 \operatorname{Bias}(G_3)=-u_q-g,\quad
 \operatorname{Bias}(G_4)=0.
\]
对应方差逐项为
\begin{align*}
 \Var(G_1)&=\E_p\lVert s_A\rVert^2,\\
 \Var(G_2)&=\E_p[r_A^2\lVert s_A\rVert^2]-\lVert g\rVert^2,\\
 \Var(G_3)&=\E_p[(1-e^{-r_A})^2\lVert s_A\rVert^2]-\lVert u_q\rVert^2,\\
 \Var(G_4)&=\E_p[r_A^2\lVert s_A\rVert^2]-\lVert g\rVert^2.
\end{align*}
例如 $\E_p[G_3]=\E_p[s_A]-\E_p[(q(A)/p(A))s_A]=-u_q$。由此得到两个明确区别：K2 的数值有偏但其固定样本反传梯度无偏；K1 与 K3 的数值无偏但其裸反传梯度通常有偏。这些期望在每次抽样均使用当前 $p$ 时成立；不能将固定样本的导数与包含采样测度导数的 $\nabla_\theta\E_{p_\theta}$ 混用。

确定性 Top-$J$ 数值目标的梯度也可以显式计算。记
\[
 d_P=\nabla_\theta P=\sum_{v\in S}p(v)s_v,
 \qquad \ell_{\tail}=\log(\tau_p/\tau_q).
\]
由乘积法则与 $\nabla_\theta\tau_p=-d_P$，有
\begin{align*}
 \nabla_\theta H&=h_g+d_P,\\
 \nabla_\theta(H+M)&=h_g-\ell_{\tail}d_P=h_g+\tau_pb_{\mathrm{TA}},\\
 \nabla_\theta K_J^{\mathrm{norm}}
 &=\frac{h_g}{P}-\frac{H}{P^2}d_P.
\end{align*}
所以三者相对于完整 KL 梯度的偏差依次为
\[
 d_P-\tau_p\mu_g,\qquad
 \tau_p(b_{\mathrm{TA}}-\mu_g),\qquad
 \frac{h_g}{P}-\frac{H}{P^2}d_P-g,
\]
方差均为零。尤其 $h_g$ 是停止系数的头部 score 梯度，$\nabla H=h_g+d_P$ 是头部数值的真实梯度；两者必须分别命名。

\subsection{停止梯度 surrogate 的显式构造}

为把数学上的梯度向量对应到可反传的标量，记 $\operatorname{sg}$ 为停止梯度算子。固定抽样动作和事件，定义
\begin{align*}
 L_{\mathrm{HETS}}
 &=\sum_{v\in S}\operatorname{sg}(p(v)r_v)\log p(v)
   +J_A\operatorname{sg}(r_A)\log p(A),\\
 L_{\mathrm{Res}}
 &=L_{\mathrm{HETS}}
   +\operatorname{sg}((\tau_p-J_A)\ell_{\tail})\log\tau_p.
\end{align*}
逐项反传，得到
\begin{align*}
 \nabla L_{\mathrm{HETS}}&=h_g+J_Ar_As_A,\\
 \nabla L_{\mathrm{Res}}
 &=h_g+J_Ar_As_A+(\tau_p-J_A)\ell_{\tail}\nabla\log\tau_p\\
 &=h_g+\tau_pb_{\mathrm{TA}}+J_A(r_As_A-b_{\mathrm{TA}}).
\end{align*}
这两条等式依赖停止整个系数的梯度，包括其中的 $p(v)$、$r_v$、$\tau_p$ 和 $\ell_{\tail}$；若遗漏停止操作，就会引入额外导数。

若希望 forward 展示某个数值估计量 $F$，同时 backward 使用 $L$ 的梯度，可定义
\[
 \widetilde L=\operatorname{sg}(F)+L-\operatorname{sg}(L).
\]
逐样本有 $\widetilde L=F$，但 $\nabla\widetilde L=\nabla L$。因此 forward 的数值偏差和 backward 的梯度偏差可以独立设计，审计时必须同时给出 $F$ 和 $L$。

\subsection{K4 尾部反传与 HETS 的精确关系}
\label{sec:kl-k4-hets}

前文 HETS 数值估计量在项目中也称为 EMA，其公式为 $\widehat K^{\mathrm{HETS}}=H+J_Ar_A$。若头部按真实数值 $H$ 求导，而尾部使用本章明确命名的 K4$^{\mathrm{live}}$ 构造，则自动微分得到
\[
 \widehat g^{\mathrm{EMA\text{-}K4}}
 =h_g+d_P+J_A(r_A+1)s_A,\qquad d_P=\nabla_\theta P.
\]
这与 HETS 的 $h_g+J_Ar_As_A$ 相差 $d_P+J_As_A$。因为 $\E_p[J_As_A]=\nabla_\theta\tau_p=-d_P$，差值期望为零，但样本和方差均不同。

\begin{theorem}[头尾 score 常数的匹配]
对固定实数 $\alpha$，定义
\[
 \widehat g_\alpha
 =h_g+\alpha d_P+J_A(r_A+\alpha)s_A.
\]
则它对 $g$ 无偏，且
\begin{equation}
 \Var(\widehat g_\alpha)
 =\sum_{v\notin S}p(v)(r_v+\alpha)^2\lVert s_v\rVert^2
  -\lVert\tau_p\mu_g-\alpha d_P\rVert^2.
 \label{eq:kl-alpha-var}
\end{equation}
$\alpha=0$ 是 HETS，$\alpha=1$ 是真实头部数值加 masked K4$^{\mathrm{live}}$ 尾部的反传。
\end{theorem}

\begin{proof}
期望为
\[
 h_g+\alpha d_P+\tau_p\mu_g+\alpha\sum_{v\notin S}p(v)s_v
 =h_g+\tau_p\mu_g=g.
\]
确定性头部不贡献协方差；随机尾部的均值为 $\tau_p\mu_g-\alpha d_P$，二阶矩为式~\eqref{eq:kl-alpha-var} 的求和项，两者相减得到方差。
\end{proof}

如果误把真实头部梯度 $h_g+d_P$ 与尾部的 $J_Ar_As_A$ 拼接，得到的期望为 $g+d_P$，偏差为 $d_P$。若头部停止系数为 $h_g$，尾部却使用 $J_A(r_A+1)s_A$，偏差为 $-d_P$。两端的 score 常数必须匹配；仅凭“精确头部加 sampled tail”这句话无法判断无偏性。

EMA-PG 固定版本的 \texttt{compute\_memory\_efficient\_kl} 在 \texttt{full\_reverse} 路径中直接反传头部 $H$，尾部使用停止 log-ratio 的 K4。因此在共同固定支持、无 clipping 且尾部测度校正正确时，forward 无偏，backward 的偏差为 $d_P$。actor 传入的 log-importance weight 已 detach，不能通过它的导数消掉 $d_P$。\footnote{同一固定版本的 actor 路径见 \href{https://github.com/LunjunZhang/ema-pg/blob/f9b977632ad948ce935432a517c22291bfaba562/search/verl/workers/actor/dp_actor.py}{\texttt{dp\_actor.py}}。}本章所证明的无偏 EMA 数值结构与论文 HETS 的向量结构，不能据名称相近就认定为该路径的逐样本反传行为。

本仓库的 \path{tests/test_kl_estimators.py} 核验单 token K4$^{\mathrm{live}}$ 的数值与梯度，对应 $(r_A+1)s_A$。\path{experiments/analysis/ema_topk_k_analysis.py} 中的结构统计采用 $h_g+d_P$ 头部与 $r_As_A$ 尾部，其期望具有上述 $d_P$ 偏差；这些统计不能作为正确匹配头尾常数的估计量或 HETS 的无偏性证据。\path{src/mirroropd/training/losses.py} 的 OPD baseline 使用停止梯度的 Tail-aware 标量，属于 scalar baseline 路径。本章的数学分类分别描述这三个实际对象。
